%! Author = angel

% Preamble
\documentclass[11pt]{article}

% Packages
\usepackage{amsmath,amssymb,amsthm}
\usepackage[letterpaper,margin=0.75in]{geometry}
\usepackage{pgfplots}
\usepackage{gensymb}
\usepackage{tikz}
\usepackage[usenames, dvipsnames]{xcolor}
\usepackage{fancyhdr}
\pgfplotsset{compat=1.18}
\usepackage{tikz-cd}

\usetikzlibrary{decorations.markings, arrows.meta}
% Document
\tikzset{
    marrow/.style={decoration={markings,mark=at position 0.5 with {\arrow{#1}}}, postaction=decorate}
}


\newcommand{\fahrenheit}{\degree \text{F}}

% Document
\begin{document}
    \noindent \textbf{Chapter 37: Relativity}

    \vspace{1em}
    \noindent Space and time are entangled,
    in which the time between two events rely on the distance between the two.
    Events are generally defined as something that happens, where we can assign three space coordinates ($x,y,z$), one time coordinate.
    Inertial reference frames are the perspectives we take when observing an object that is at rest or moving with constant speed.
    An observer, in an inertial reference frame, can assign an event spacetime coordinates since space and time are entangled.\\
    
    \noindent \vspace{1em}
    Relativity measures how space and time are entangled, following the two postulates:
    \begin{enumerate}
      \item \textbf{The Relativity Postulate:}
        The laws of Physics apply to all observers regardless of inertial reference frame.
      \item \textbf{The Speed of Light Postulate:}
        The speed of light in  a vacuum will always be $c$ regardless of direction and inertial reference frame.
    \end{enumerate}    
    We define the speed of light $c$ (m/s) as:
    \begin{equation}
      c = 299,792,458 \text{ m/s } = 3.00 \times 10^8 \text{ m/s} \tag{speed of light}
    \end{equation}
    We can also call this \textit{ultimate speed} as the more kinetic energy speed a particle has, the closer it will get to this boundary.
    \vspace{1em}

    \noindent Assigning spacetime coordinates to an event allows us to see how time is relative.
    Let's say there was a firework exploding in the sky a mile away, and you turn on your phone both at exactly 8:00 PM.
    From your own observation, you notice your phone turn on first, then the firework expolsion.
    As the light from the firework has to travel longer than the phone in your hand,
    these two events seem to happen at different types from your inertial frame of reference.
    
    \vspace{1em}
    \noindent Now, lets take a look as an observer on a plane traveling at the exact midpoint of the two events.
    In this case, the observer notices the firework and the orignial observer's phone at the same time,
    as the distances between the observer and distances changed.
    The new observer sees the event happen \textit{simultaneously}.
    We find that simultaneity is relative, not absolute.
\end{document}
